\documentclass[aspectratio=169,11pt]{beamer}
\usetheme{Madrid}
\usecolortheme{dolphin}
\setbeamertemplate{navigation symbols}{}
\setbeamertemplate{caption}[numbered]

\usepackage[utf8]{inputenc}
\usepackage[T1]{fontenc}
\usepackage[spanish,es-noshorthands]{babel}
\usepackage{amsmath,amssymb,amsthm,mathtools}
\usepackage{tikz}
\usepackage{pgfplots}
\pgfplotsset{compat=1.18}
\usepackage{booktabs}
\usepackage{array}

\title[Prob. y Estadística]{Clase 22: Métodos de Estimación}
\subtitle{Unidad 6: Estimación de parámetros}
\institute[UNS]{Universidad Nacional del Sur \\ Departamento de Matemática}
\author{Probabilidad y Estadística (5765)}
\date{}

\begin{document}

\begin{frame}
\titlepage
\end{frame}

\begin{frame}{Contenidos}
\tableofcontents
\end{frame}

\section{Método de mínimos cuadrados}

\begin{frame}{Idea general del método de mínimos cuadrados}
\begin{block}{Motivación}
Dado un conjunto de observaciones $x_1,\dots,x_n$ que en teoría deberían ajustarse a un modelo con parámetro(s) $\theta$, el \textbf{método de mínimos cuadrados (MC)} elige como estimador el valor $\widehat\theta$ que \textbf{minimiza la suma de los cuadrados de los errores} entre los datos observados y los valores predichos por el modelo.
\end{block}

\begin{definition}[Estimador de mínimos cuadrados]
Si el modelo predice $E(X_i) = g(\theta)$, el estimador de MC de $\theta$ es el que minimiza
\[
Q(\theta) = \sum_{i=1}^n \left(x_i - g(\theta)\right)^2.
\]
\end{definition}

\begin{alertblock}{Uso típico}
Es la base de la \textbf{regresión lineal} (Unidad 8): allí $g$ depende de variables explicativas y se minimiza la suma de residuos al cuadrado. Aquí lo presentamos en su versión más simple, para fijar la idea.
\end{alertblock}
\end{frame}

\begin{frame}{Ejemplo: mínimos cuadrados para estimar $\mu$}
\begin{example}
Sea $X_1,\dots,X_n$ m.a.\ con $E(X_i)=\mu$ (modelo $g(\mu)=\mu$ constante). Buscamos $\widehat\mu$ que minimice
\[
Q(\mu) = \sum_{i=1}^n (x_i-\mu)^2.
\]
Derivando e igualando a cero:
\[
Q'(\mu) = -2\sum_{i=1}^n (x_i-\mu) = 0 \;\Longrightarrow\; \sum_{i=1}^n x_i = n\mu \;\Longrightarrow\; \widehat\mu = \frac{1}{n}\sum_{i=1}^n x_i = \overline x.
\]
Como $Q''(\mu) = 2n>0$, es efectivamente un mínimo. \textbf{Conclusión:} el estimador de mínimos cuadrados de $\mu$ coincide con la media muestral $\overline X$.
\end{example}
\end{frame}

\section{Función de verosimilitud}

\begin{frame}{Motivación del método de máxima verosimilitud}
\begin{block}{Pregunta central}
Dada una muestra observada $x_1,\dots,x_n$, ¿qué valor de $\theta$ hace que \textbf{sea más "verosímil" (plausible)} haber observado exactamente esos datos?
\end{block}

\begin{definition}[Función de verosimilitud]
Sea $X_1,\dots,X_n$ m.a.\ de una población con función de probabilidad/densidad $f(x;\theta)$. La \textbf{función de verosimilitud} es
\[
L(\theta) = L(\theta \mid x_1,\dots,x_n) = \prod_{i=1}^n f(x_i;\theta),
\]
vista como función de $\theta$ (los datos $x_1,\dots,x_n$ están fijos, son los observados).
\end{definition}

\begin{alertblock}{Diferencia clave con la densidad conjunta}
$f(x_1,\dots,x_n;\theta)$ es función de los datos, con $\theta$ fijo. $L(\theta)$ es la \emph{misma expresión} pero pensada como función de $\theta$, con los datos fijos.
\end{alertblock}
\end{frame}

\begin{frame}{Estimador de máxima verosimilitud (EMV)}
\begin{definition}[EMV]
El \textbf{estimador de máxima verosimilitud} de $\theta$ es el valor $\widehat\theta$ que maximiza $L(\theta)$:
\[
\widehat\theta_{EMV} = \operatorname*{arg\,max}_{\theta \in \Theta} L(\theta).
\]
\end{definition}

\begin{block}{Log-verosimilitud}
Como $\log$ es una función estrictamente creciente, maximizar $L(\theta)$ equivale a maximizar
\[
\ell(\theta) = \ln L(\theta) = \sum_{i=1}^n \ln f(x_i;\theta),
\]
que en general es mucho más fácil de derivar (transforma productos en sumas).
\end{block}

\begin{alertblock}{Receta general}
1) Escribir $L(\theta)$. \quad 2) Tomar $\ell(\theta)=\ln L(\theta)$. \quad 3) Resolver $\dfrac{d\ell}{d\theta}=0$. \quad 4) Verificar que es máximo ($\ell''<0$).
\end{alertblock}
\end{frame}

\section{EMV en distribuciones discretas}

\begin{frame}{EMV para la distribución Bernoulli}
\begin{example}
Sea $X_1,\dots,X_n$ m.a.\ de $X\sim \text{Be}(p)$, con $f(x;p) = p^x(1-p)^{1-x}$, $x\in\{0,1\}$.
\[
L(p) = \prod_{i=1}^n p^{x_i}(1-p)^{1-x_i} = p^{\sum x_i}(1-p)^{n-\sum x_i}.
\]
Sea $t=\sum_{i=1}^n x_i$ (cantidad de éxitos). Entonces
\[
\ell(p) = t\ln p + (n-t)\ln(1-p).
\]
Derivando:
\[
\ell'(p) = \frac{t}{p} - \frac{n-t}{1-p} = 0 \;\Longrightarrow\; t(1-p) = (n-t)p \;\Longrightarrow\; t = np \;\Longrightarrow\; \widehat p = \frac{t}{n} = \overline x.
\]
Como $\ell''(p) = -t/p^2 - (n-t)/(1-p)^2 < 0$, es máximo. \textbf{El EMV de $p$ es la proporción muestral} $\widehat p = \overline X$.
\end{example}
\end{frame}

\begin{frame}{EMV para la distribución de Poisson}
\begin{example}
Sea $X_1,\dots,X_n$ m.a.\ de $X\sim \text{Poisson}(\lambda)$, $f(x;\lambda)=\dfrac{e^{-\lambda}\lambda^x}{x!}$.
\[
L(\lambda) = \prod_{i=1}^n \frac{e^{-\lambda}\lambda^{x_i}}{x_i!} = \frac{e^{-n\lambda}\lambda^{\sum x_i}}{\prod x_i!}.
\]
\[
\ell(\lambda) = -n\lambda + \left(\sum_{i=1}^n x_i\right)\ln\lambda - \ln\left(\prod x_i!\right).
\]
\[
\ell'(\lambda) = -n + \frac{\sum x_i}{\lambda} = 0 \;\Longrightarrow\; \widehat\lambda = \frac{1}{n}\sum_{i=1}^n x_i = \overline x.
\]
Como $\ell''(\lambda) = -\sum x_i/\lambda^2 < 0$ (si algún $x_i>0$), es máximo. \textbf{El EMV de $\lambda$ es $\overline X$.}
\end{example}
\end{frame}

\section{EMV en distribuciones continuas}

\begin{frame}{EMV para la distribución exponencial}
\begin{example}
Sea $X_1,\dots,X_n$ m.a.\ de $X\sim \text{Exp}(\lambda)$, $f(x;\lambda)=\lambda e^{-\lambda x}$, $x>0$.
\[
L(\lambda) = \prod_{i=1}^n \lambda e^{-\lambda x_i} = \lambda^n e^{-\lambda \sum x_i}.
\]
\[
\ell(\lambda) = n\ln\lambda - \lambda\sum_{i=1}^n x_i.
\]
\[
\ell'(\lambda) = \frac{n}{\lambda} - \sum_{i=1}^n x_i = 0 \;\Longrightarrow\; \widehat\lambda = \frac{n}{\sum_{i=1}^n x_i} = \frac{1}{\overline x}.
\]
Como $\ell''(\lambda)=-n/\lambda^2<0$, es máximo. \textbf{El EMV de $\lambda$ es $1/\overline X$} (nótese: $\widehat\lambda$ es sesgado, aunque consistente, ya que $E(1/\overline X)\neq 1/E(\overline X)$ en general).
\end{example}
\end{frame}

\begin{frame}{EMV para la distribución normal (deducción completa)}
\begin{example}
Sea $X_1,\dots,X_n$ m.a.\ de $X\sim N(\mu,\sigma^2)$, ambos parámetros desconocidos.
\[
f(x;\mu,\sigma^2) = \frac{1}{\sqrt{2\pi\sigma^2}}\exp\left(-\frac{(x-\mu)^2}{2\sigma^2}\right).
\]
\[
L(\mu,\sigma^2) = (2\pi\sigma^2)^{-n/2}\exp\left(-\frac{1}{2\sigma^2}\sum_{i=1}^n (x_i-\mu)^2\right).
\]
\[
\ell(\mu,\sigma^2) = -\frac{n}{2}\ln(2\pi) - \frac{n}{2}\ln(\sigma^2) - \frac{1}{2\sigma^2}\sum_{i=1}^n (x_i-\mu)^2.
\]
\end{example}
\end{frame}

\begin{frame}{EMV para la normal: resolución del sistema}
Derivando respecto de $\mu$ e igualando a cero:
\[
\frac{\partial \ell}{\partial \mu} = \frac{1}{\sigma^2}\sum_{i=1}^n (x_i-\mu) = 0 \;\Longrightarrow\; \widehat\mu = \frac{1}{n}\sum_{i=1}^n x_i = \overline x.
\]
Derivando respecto de $\sigma^2$ (tratándolo como una única variable) e igualando a cero:
\[
\frac{\partial \ell}{\partial \sigma^2} = -\frac{n}{2\sigma^2} + \frac{1}{2\sigma^4}\sum_{i=1}^n (x_i-\mu)^2 = 0
\]
\[
\Longrightarrow \; \widehat\sigma^2 = \frac{1}{n}\sum_{i=1}^n (x_i-\widehat\mu)^2 = \frac{1}{n}\sum_{i=1}^n (x_i-\overline x)^2.
\]
\begin{alertblock}{Observación importante}
El EMV de $\sigma^2$ es $\widehat\sigma^2$ (con $n$, \textbf{no} $n-1$): coincide con el estimador \textbf{sesgado} de la Clase 21, no con $S^2$. El EMV no siempre es insesgado.
\end{alertblock}
\end{frame}

\section{Propiedades del EMV}

\begin{frame}{Propiedades del estimador de máxima verosimilitud}
\begin{block}{Propiedad de invarianza}
Si $\widehat\theta$ es el EMV de $\theta$ y $g$ es una función biunívoca (inyectiva), entonces $g(\widehat\theta)$ es el EMV de $g(\theta)$.
\end{block}
\begin{example}
Como el EMV de $\sigma^2$ es $\widehat\sigma^2 = \frac{1}{n}\sum(x_i-\overline x)^2$, por invarianza el EMV del desvío estándar $\sigma$ es
\[
\widehat\sigma = \sqrt{\widehat\sigma^2} = \sqrt{\frac{1}{n}\sum_{i=1}^n (x_i-\overline x)^2}.
\]
\end{example}

\begin{block}{Propiedades asintóticas (bajo condiciones de regularidad)}
\begin{itemize}
\item \textbf{Consistencia:} $\widehat\theta_{EMV} \xrightarrow{P} \theta$ cuando $n\to\infty$.
\item \textbf{Eficiencia asintótica:} su varianza se aproxima a la cota de Cramér--Rao para $n$ grande.
\item \textbf{Normalidad asintótica:} $\widehat\theta_{EMV}$ tiene distribución aproximadamente normal para $n$ grande.
\end{itemize}
\end{block}
\end{frame}

\section{Ejemplo numérico integrador}

\begin{frame}{Ejemplo numérico: EMV con datos concretos}
\begin{example}
Se registra el número de llamados que recibe un call center por minuto, durante $n=8$ minutos, asumiendo $X\sim\text{Poisson}(\lambda)$:
\[
3,\ 5,\ 2,\ 4,\ 3,\ 6,\ 4,\ 5.
\]
Sabemos que el EMV de $\lambda$ en Poisson es $\widehat\lambda = \overline x$. Calculamos:
\[
\sum_{i=1}^8 x_i = 3+5+2+4+3+6+4+5 = 32, \qquad \widehat\lambda = \overline x = \frac{32}{8} = 4.
\]
\textbf{Estimación de máxima verosimilitud:} $\widehat\lambda = 4$ llamados por minuto.

Si además nos interesara $g(\lambda) = P(X=0) = e^{-\lambda}$ (probabilidad de no recibir llamados en un minuto), por invarianza:
\[
\widehat{g} = e^{-\widehat\lambda} = e^{-4} \approx 0.0183.
\]
\end{example}
\end{frame}

\section{Resumen}

\begin{frame}{Resumen}
\begin{block}{Método de mínimos cuadrados}
Minimiza $\sum (x_i - g(\theta))^2$. Para el modelo de media constante, coincide con $\overline X$.
\end{block}

\begin{block}{Método de máxima verosimilitud}
Maximiza $L(\theta)=\prod f(x_i;\theta)$, o equivalentemente $\ell(\theta)=\ln L(\theta)$.
\end{block}

\begin{center}
\begin{tabular}{lcc}
\toprule
\textbf{Distribución} & \textbf{Parámetro} & \textbf{EMV} \\
\midrule
Bernoulli($p$) & $p$ & $\widehat p = \overline X$ \\
Poisson($\lambda$) & $\lambda$ & $\widehat\lambda = \overline X$ \\
Exponencial($\lambda$) & $\lambda$ & $\widehat\lambda = 1/\overline X$ \\
Normal($\mu,\sigma^2$) & $\mu$ & $\widehat\mu = \overline X$ \\
Normal($\mu,\sigma^2$) & $\sigma^2$ & $\widehat\sigma^2 = \frac{1}{n}\sum(X_i-\overline X)^2$ \\
\bottomrule
\end{tabular}
\end{center}

\begin{alertblock}{Próxima clase}
Usaremos estos estimadores puntuales como base para construir \textbf{intervalos de confianza}: en vez de dar un único valor, daremos un rango de valores plausibles para $\theta$.
\end{alertblock}
\end{frame}

\end{document}
